\section{Introduction}\label{sec:introduction}

Which coordinates of a decision problem can be hidden without changing the optimal action, and how hard is it to certify that claim exactly? We expected one answer. We found three fundamentally different ones. We study the same underlying question across static, stochastic, and sequential regimes. For a decision problem $\mathcal{D}=(A,S,U)$ with $S = X_1 \times \cdots \times X_n$, a coordinate set $I$ is sufficient when agreement on $I$ forces agreement on the optimal-action set:
\[
s_I = s'_I \implies \Opt(s) = \Opt(s').
\]
This is the basic problem of exact relevance certification: how much of the state must remain visible to preserve the decision boundary, and how hard it is to certify that the rest is irrelevant. The question resembles subset selection, but exact certification has a different quantifier structure and, as the paper shows, a different complexity landscape.

This is a certification problem rather than a forward-evaluation problem. Evaluating $U(a,s)$ or $\Opt(s)$ on a fully specified state is one task; proving that a family of coordinates is irrelevant requires ruling out every counterexample pair of states. That universal structure drives the static regime. A first structural point is worth stating up front: in the static regime, a set is sufficient exactly when it contains every relevant coordinate. So \textsc{Minimum-Sufficient-Set} is not a genuinely alternating search problem after all; the minimum sufficient set is just the relevant-coordinate set. At the same time, the optimizer quotient is not an arbitrary definition: it is the canonical coarsest abstraction that preserves optimal-action distinctions. \textsc{Anchor-Sufficiency} does not collapse in the same way, because the anchor assignment is itself part of the existential choice and must then satisfy a universal fiberwise condition. At this static boundary the paper touches classical reduct-style attribute reduction, but its main object is broader: the same underlying question is then transported to stochastic conditioning and to sequential models with temporal structure.

\paragraph{Why probability changes the exact question.}
In stochastic decision problems, exact analysis separates into a preservation question and a decisiveness question. Preservation asks whether coarse information reproduces the full-information optimizer: does the optimizer induced by the retained coordinates agree with the optimizer seen when the entire state is known? Decisiveness asks whether the coarse information itself yields a unique Bayes-optimal action on each observable fiber: is the retained signal itself a complete decision interface? Formally, these are \emph{stochastic sufficiency} and \emph{stochastic decisiveness}. The first is a fidelity notion; the second is a decision-completeness notion. They are distinct, and probability makes both natural. Preservation gives the direct stochastic analogue of static sufficiency and the bridge back to the optimizer quotient; under full support, its minimum and anchor variants inherit the static coNP and $\Sigma_2^P$ classifications. Decisiveness yields the strongest succinct-encoding complexity classification established in the paper.

This is also what the complexity classification buys beyond nearby literatures on reducts, feature selection, or abstraction heuristics. Those traditions often ask whether some smaller representation is informative enough, predictive enough, or value-preserving enough. The present paper asks a sharper question: when can one certify exactly that hiding coordinates leaves the optimal-action correspondence unchanged? That exactness requirement is what later yields a no-general exact minimizer theorem, a regime-sensitive tractability boundary, and a simple but important reliability corollary about exact certifiers in the hard regime. The point is not merely that some small interface may work in practice, but whether there is a certificate that no omitted coordinate can change the decision.

This distinction also explains why exactness matters. When a system replaces exact relevance certification with a heuristic simplification, the omitted relevance does not disappear; it is shifted into residual uncertainty that must be handled elsewhere in the system. We refer to that externalized burden as the \emph{simplicity tax}. The results of this paper show that this tax is not merely rhetorical: exact simplification can be computationally expensive, but inexact simplification does not erase decision-relevant structure; it relocates it.

\paragraph{Informal interpretive guide.}
The formal objects studied here admit a compact intuitive reading. A relevant coordinate is a point of \emph{leverage}: changing it can change the decision. The optimizer quotient isolates the decision \emph{signal}: it is the coarsest abstraction that still preserves optimal-action distinctions. The static collapse expresses \emph{inevitability}: in the static regime, the optimal simplification is forced by the structure of relevance rather than discovered by search. The simplicity tax is \emph{liability}: omitted exact relevance does not disappear, but must be handled elsewhere in the system. At the task level, sufficiency is \emph{observation}, anchor sufficiency is \emph{discovery}, and sequential certification is \emph{exposure}: omitted information may appear harmless now while still creating future liability elsewhere in the system.

The paper's central claim is that this is not a collection of isolated hardness facts. Once the same underlying question is tracked across richer input models, the difficulty escalates in a systematic way. Static certification is governed by counterexample exclusion. In the stochastic regime, preservation yields explicit-state tractability, bridges to the static theory, and full-support inheritance for the minimum and anchor variants, while decisiveness yields the strongest succinct-encoding complexity classification. Sequential certification is governed by temporally induced optimizer structure. Read as a regime-sensitive complexity matrix, adding probability introduces conditional-comparison hardness on the decisiveness side, and adding time lifts the core queries again to PSPACE. The optimizer quotient packages the same issue structurally: it is the coarsest abstraction of the state space that preserves optimal-action distinctions.

Throughout, we work with finite coordinate-structured decision problems under the encoding conventions fixed in Section~\ref{sec:encoding}.

\subsection{Contributions}

The paper is best read as one theorem package with five linked takeaways: a structural static characterization of minimum sufficiency, a regime-sensitive complexity matrix, an encoding-sensitive tractability contrast, downstream corollaries for exact certifiers in the hard regime, and a mechanically checked finite-decision core.

\begin{enumerate}
\item \textbf{Static Structural Theorem.} In the static regime, sufficient sets are exactly the supersets of the relevant-coordinate set. This structural characterization identifies the minimum sufficient set directly and yields coNP-completeness of \textsc{Minimum-Sufficient-Set} as a consequence. By contrast, \textsc{Anchor-Sufficiency} remains $\Sigma_2^P$-complete because its existential anchor choice does not collapse.

\item \textbf{Stochastic Preservation and Decisiveness.} We track the same underlying question across static, stochastic, and sequential regimes. The static regime has a complete collapse-and-classification theorem package. In the stochastic regime, stochastic sufficiency is polynomial-time under explicit-state encoding, mechanically bridged back to static sufficiency and the optimizer quotient, and its minimum and anchor variants inherit coNP-completeness and $\Sigma_2^P$-completeness under full support. Stochastic decisiveness is polynomial-time under explicit-state encoding and PP-hard under succinct encoding. We retain the traditional names ``stochastic anchor sufficiency'' and ``stochastic minimum sufficiency'' for the anchor and minimum variants of the decisiveness predicate; they are PP-hard and lie in $\textsf{NP}^{\textsf{PP}}$, with those oracle-class memberships proved in the paper text and their finite witness/checking cores mechanized in the artifact. The sequential row gives the PSPACE end of the story.

\item \textbf{Encoding-sensitive tractability contrast.} We prove an explicit-state versus succinct contrast: logarithmic-size sufficient sets admit polynomial-time algorithms in the explicit regime, while linear-size sufficient sets in the succinct regime inherit ETH-conditioned exponential lower bounds. We also isolate six structured tractable subcases under standard restrictions.

\item \textbf{Methodological Consequences for Exact Certification.} In the hard regime, exact certifiers cannot simultaneously be sound, complete, and polynomial-budgeted. This appears in the paper as a direct complexity-theoretic corollary rather than as a separate structural theorem. The same section also isolates witness-checking lower bounds, approximation obstructions, and the externalization of omitted exact relevance.

\item \textbf{Mechanized theorem infrastructure.} The main reductions, hardness packages, finite deciders, search procedures, and explicit-state step-counting wrappers are mechanically checked \leanmeta{\LHrng{DC}{91}{99}, \LHrng{OU}{1}{2}, \LHrng{OU}{8}{12}, \LHrng{EH}{1}{10}.} The artifact mechanizes the finite witness/checking schemata and reduction cores used by the stochastic upper-bound arguments. The corresponding oracle-class memberships are proved in the paper text, following standard complexity-theoretic proof conventions. Thus the formalization provides independent verification of the finite combinatorial core underlying those results.
\end{enumerate}

\subsection{Quantifier Structure}

The complexity gap comes from the structure of the question. Insufficiency has short witnesses: two states that agree on the proposed coordinates but induce different optimal-action sets. Sufficiency, by contrast, asks for the absence of every such witness. So the object-level task is forward evaluation on a fully specified state, while the meta-level task is universal verification over counterexample pairs. This difference drives the coNP classification and explains why relevance identification is not just another forward-evaluation problem.

The static regime also contains the paper's main structural simplification. Although minimum sufficiency can be written with an outer existential quantifier over coordinate sets, the existential layer collapses once sufficiency is characterized as containment of all relevant coordinates. The stochastic and sequential regimes do not simplify in the same way: the stochastic regime separates into preservation and decisiveness, with full-support inheritance on the preservation side and counting-style hardness on the decisiveness side, while temporal evolution introduces dependence on temporally induced optimizer structure.

\subsection{Paper Structure}

Section~\ref{sec:formal-setup} introduces the formal setup and encoding conventions. Sections~\ref{sec:hardness}, \ref{sec:stochastic}, \ref{sec:sequential}, and~\ref{sec:regime-hierarchy} develop the core theorem package. Sections~\ref{sec:dichotomy} and~\ref{sec:tractable} give the encoding-sensitive contrast and tractable subcases. Section~\ref{sec:engineering-corollaries} collects implications for exact simplification and reliability consequences for exact certifiers. Section~\ref{sec:related} situates the work relative to formalized complexity and decision theory.

\subsection{Artifact Availability}

The proof artifact is archived at \url{https://doi.org/10.5281/zenodo.18140966}. The cited-content snapshot contains \LeanReleaseLines\ lines across \LeanReleaseFiles\ files with \LeanReleaseSorry\ occurrences of \texttt{sorry}.
